% §13.1 Algebraic Types Recap
%
% Monograph register: the four type constructors, each prior structure as a
% Bernoulli[T]. Fix: Bernoulli preserves product/list/function (Kronecker), but
% the SUM is the exception (the tag error couples the components) -- per the source
% (composite_types.tex) and Core Principle 2. Cross-refs ch:approximate-sets,
% ch:set-composition, ch:approximate-functions, ch:sketches-estimators,
% ch:approximate-relations, sec:bernoulli-of-composite, sec:composition-functor.
% DEFINED: sec:algebraic-types-recap.

\section{Algebraic Types Recap}
\label{sec:algebraic-types-recap}

Twelve chapters in, the same calculations keep recurring. The Kronecker product
of channel matrices that organized $m$ Boolean queries in chapter~4 returned for
the n-ary channels of chapter~9 and the relational products of chapter~12; the
set-composition rules of chapter~6 lined up with chapter~9's; CountMin and
HyperLogLog were noisy counters, a different output type under the same two
axioms. The thread is that each structure is a $\Bernoulli{T}$ for some algebraic
type $T$\index{approximate type}. This chapter organizes the structures by type.
The organization has direct predecessors in approximate and probabilistic
programming, where type qualifiers carry a correctness probability through a
program (see the Notes and Lineage at the end of the chapter); what follows records which constructors
make the error law of a composite value factor, under which independence
hypotheses, and where factorization fails.

\paragraph{The four constructors.}
Algebraic data types are built from four constructors: the
\emph{product}\index{product type} $T \times U$ (pairs, records, structs;
cardinality $|T||U|$), the \emph{sum}\index{sum type} $T + U$ (tagged unions,
variants, enums; cardinality $|T| + |U|$, with a tag recording which summand),
the \emph{list}\index{recursive type} $L T$ (finite sequences of any length),
and the \emph{function} $T \to U$ (the most general, into which the
others embed). They close under composition, so any type built from primitives
and these four has a clear structural decomposition.

Each approximate structure so far lives at one constructor. An
\emph{approximate set}\index{approximate set} (\Cref{ch:approximate-sets}) is a noisy predicate
$\Bernoulli{T \to \BoolType}$, with the chapter-6 composition rules
(\Cref{ch:set-composition}) the rules for composing such predicates. An
\emph{approximate map or channel} (\Cref{ch:approximate-functions}) is a
$\Bernoulli{T \to U}$ for finite $T, U$, the n-ary channel. A \emph{sketch}
(\Cref{ch:sketches-estimators}) is a noisy value of an aggregate type, CountMin
and HyperLogLog as $\Bernoulli{\mathbb{N}}$, MinHash as $\Bernoulli{[0,1]}$. An
\emph{approximate relation}\index{approximate relation}
(\Cref{ch:approximate-relations}) is a noisy
predicate on a product, $\Bernoulli{T \times U \to \BoolType}$. The type tells
how the joint error law of a composite value factors when the components' noise
is independent; it does not by itself give the error of a computation on such
values (\Cref{sec:composition-functor}).

\paragraph{The lift.}
Write $\Bernoulli{-}$ as a type-level operation: $\Bernoulli{T}$ is the type of
random approximate values over $T$, with the atom $\Bernoulli{\BoolType}$ the
noisy bit of Part~I. The substantive content of the chapter is how
$\Bernoulli{-}$ behaves on each constructor.
\Cref{sec:bernoulli-of-composite} shows it factors through three of the four
when the relevant noise is independent: the product (a tensor of factor lifts),
the list (a per-position lift), and the function (a per-input lift), each by the
Kronecker factorization of chapters~4 and~9. The sum is the exception: a sum
value's \emph{tag} can itself err, coupling the two components, and there the
factorization breaks even with independent noise.
\Cref{sec:composition-functor} collects the rules, states what they do and do
not give, and identifies the sum's tag error with the membership error of
approximate sets.
